
\documentclass{article}
\usepackage{colortbl}
\usepackage{makecell}
\usepackage{multirow}
\usepackage{supertabular}

\begin{document}

\newcounter{utterance}

\centering \large Interaction Transcript for game `dond', experiment `coop', episode 0 with claude{-}opus{-}4{-}8{-}azure.
\vspace{24pt}

{ \footnotesize  \setcounter{utterance}{1}
\setlength{\tabcolsep}{0pt}
\begin{supertabular}{c@{$\;$}|p{.15\linewidth}@{}p{.15\linewidth}p{.15\linewidth}p{.15\linewidth}p{.15\linewidth}p{.15\linewidth}}
    \# & $\;$A & \multicolumn{4}{c}{Game Master} & $\;\:$B\\
    \hline

    \theutterance \stepcounter{utterance}  
    & & \multicolumn{4}{p{0.6\linewidth}}{
        \cellcolor[rgb]{0.9,0.9,0.9}{
            \makecell[{{p{\linewidth}}}]{
                \texttt{\tiny{[P1$\langle$GM]}}
                \texttt{Jūs spēlējat kooperatīvu sarunu spēli, kurā jums jāvienojas ar citu spēlētāju par to, kā sadalīt objektu kopumu starp jums.} \\
\\ 
\texttt{Noteikumi:} \\
\texttt{(a) Jūs un otrs spēlētājs saņemat objektu kopumu. Katrs no jums saņem arī noslēptu vērtības funkciju, kas norāda, cik jūs vērtējat katru objekta veidu.} \\
\texttt{(b) Jūs apmaiņājaties ziņojumus ar citu spēlētāju, lai vienotos par to, kurš saņem kurus objektus. Katrs var nosūtīt ne vairāk kā 5 ziņojumus vai beigt spēli agri, jebkurā laikā izdarot noslēptu priekšlikumu.} \\
\texttt{(c) Jums lūdz iesniedzt noslēptu priekšlikumu, norādot objektus, kurus vēlaties, formatētu kvadrātiekavās šādi: "{[}Priekšlikums: <skaitlis> <objekta nosaukums>, <skaitlis> <objekta nosaukums>, <...>{]}"} \\
\texttt{(d) Ja jūsu priekšlikumi ir papildinoši, tas ir, ir pietiekami objekti, lai apmierinātu abus priekšlikumus, katrs spēlētājs saņem reitingu, pamatojoties uz saņemto objektu vērtību summu. Pretējā gadījumā jūs abi saņemat nulle punktu.} \\
\texttt{(e) Jūsu mērķis ir maksimizēt sava reitinga un otra spēlētāja reitinga summu.} \\
\\ 
\texttt{Sāksim.} \\
\\ 
\texttt{Pieejamo objektu kopums ir:} \\
\\ 
\texttt{1 bultskrūve, 1 soma, 2 knaibles, 2 poti} \\
\\ 
\texttt{Jūsu noslēptā vērtības funkcija:} \\
\\ 
\texttt{bultskrūve: 2, soma: 2, knaibles: 0, pots: 3} \\
\\ 
\texttt{SVARĪGI: Jūsu ziņojumi, ja vien tie nav noslēpti priekšlikumi, tiek pārraidīti tieši otram spēlētājam, tāpēc neiekļaujiet nekādu atbildi uz noteikumiem vai tekstu, kas paziņo par jūsu ziņojumu. Lai izdarītu noslēptu priekšlikumu, izmantojiet norādīto formātu. Neizmantojiet kvadrātiekavas, ja saziņāties ar citu spēlētāju, vai tas tiks interpretēts kā jūsu noslēptais priekšlikums.} \\
            }
        }
    }
    & \\ \\

    \theutterance \stepcounter{utterance}  
    & \multicolumn{4}{p{0.6\linewidth}}{
        \cellcolor[rgb]{0.9,0.9,0.9}{
            \makecell[{{p{\linewidth}}}]{
                \texttt{\tiny{[P1$\rangle$GM]}}
                \texttt{Sveiki! Lai maksimizētu mūsu kopējo rezultātu, pastāstiet, kurus objektus jūs vērtējat visaugstāk. Man knaibles ir pilnīgi bezvērtīgas (0), tāpēc labprāt atdotu tās jums. Kā jūs vērtējat bultskrūvi, somu un potus?} \\
            }
        }
    }
    & & \\ \\

    \theutterance \stepcounter{utterance}  
    & & & \multicolumn{2}{p{0.3\linewidth}}{
        \cellcolor[rgb]{0.9,0.9,0.9}{
            \makecell[{{p{\linewidth}}}]{
                \texttt{\tiny{[GM$|$GM]}}
                \texttt{continue} \\
            }
        }
    }
    & & \\ \\

    \theutterance \stepcounter{utterance}  
    & & & \multicolumn{2}{p{0.3\linewidth}}{
        \cellcolor[rgb]{0.9,0.9,0.9}{
            \makecell[{{p{\linewidth}}}]{
                \texttt{\tiny{[GM$|$GM]}}
                \texttt{Sveiki! Lai maksimizētu mūsu kopējo rezultātu, pastāstiet, kurus objektus jūs vērtējat visaugstāk. Man knaibles ir pilnīgi bezvērtīgas (0), tāpēc labprāt atdotu tās jums. Kā jūs vērtējat bultskrūvi, somu un potus?} \\
            }
        }
    }
    & & \\ \\

    \theutterance \stepcounter{utterance}  
    & & \multicolumn{4}{p{0.6\linewidth}}{
        \cellcolor[rgb]{0.9,0.9,0.9}{
            \makecell[{{p{\linewidth}}}]{
                \texttt{\tiny{[GM$\rangle$P2]}}
                \texttt{Jūs spēlējat kooperatīvu sarunu spēli, kurā jums jāvienojas ar citu spēlētāju par to, kā sadalīt objektu kopumu starp jums.} \\
\\ 
\texttt{Noteikumi:} \\
\texttt{(a) Jūs un otrs spēlētājs saņemat objektu kopumu. Katrs no jums saņem arī noslēptu vērtības funkciju, kas norāda, cik jūs vērtējat katru objekta veidu.} \\
\texttt{(b) Jūs apmaiņājaties ziņojumus ar citu spēlētāju, lai vienotos par to, kurš saņem kurus objektus. Katrs var nosūtīt ne vairāk kā 5 ziņojumus vai beigt spēli agri, jebkurā laikā izdarot noslēptu priekšlikumu.} \\
\texttt{(c) Jums lūdz iesniedzt noslēptu priekšlikumu, norādot objektus, kurus vēlaties, formatētu kvadrātiekavās šādi: "{[}Priekšlikums: <skaitlis> <objekta nosaukums>, <skaitlis> <objekta nosaukums>, <...>{]}"} \\
\texttt{(d) Ja jūsu priekšlikumi ir papildinoši, tas ir, ir pietiekami objekti, lai apmierinātu abus priekšlikumus, katrs spēlētājs saņem reitingu, pamatojoties uz saņemto objektu vērtību summu. Pretējā gadījumā jūs abi saņemat nulle punktu.} \\
\texttt{(e) Jūsu mērķis ir maksimizēt sava reitinga un otra spēlētāja reitinga summu.} \\
\\ 
\texttt{Sāksim.} \\
\\ 
\texttt{Pieejamo objektu kopums ir:} \\
\\ 
\texttt{1 bultskrūve, 1 soma, 2 knaibles, 2 poti} \\
\\ 
\texttt{Jūsu noslēptā vērtības funkcija:} \\
\\ 
\texttt{bultskrūve: 3, soma: 1, knaibles: 3, pots: 0} \\
\\ 
\texttt{SVARĪGI: Jūsu ziņojumi, ja vien tie nav noslēpti priekšlikumi, tiek pārraidīti tieši otram spēlētājam, tāpēc neiekļaujiet nekādu atbildi uz noteikumiem vai tekstu, kas paziņo par jūsu ziņojumu. Lai izdarītu noslēptu priekšlikumu, izmantojiet norādīto formātu. Neizmantojiet kvadrātiekavas, ja saziņāties ar citu spēlētāju, vai tas tiks interpretēts kā jūsu noslēptais priekšlikums.} \\
\\ 
\texttt{Sveiki! Lai maksimizētu mūsu kopējo rezultātu, pastāstiet, kurus objektus jūs vērtējat visaugstāk. Man knaibles ir pilnīgi bezvērtīgas (0), tāpēc labprāt atdotu tās jums. Kā jūs vērtējat bultskrūvi, somu un potus?} \\
            }
        }
    }
    & \\ \\

    \theutterance \stepcounter{utterance}  
    & & & \multicolumn{4}{p{0.6\linewidth}}{
        \cellcolor[rgb]{0.9,0.9,0.9}{
            \makecell[{{p{\linewidth}}}]{
                \texttt{\tiny{[GM$\langle$P2]}}
                \texttt{Sveiki! Lieliski, tas labi saskan. Man savukārt poti (pots) ir pilnīgi bezvērtīgi (0), tāpēc jūs varat paņemt abus potus. Es augstu vērtēju bultskrūvi un knaibles. Soma man ir tikai nedaudz vērtīga. Kā jums ir ar bultskrūvi un somu? Ja jūs tās nevērtējat, es labprāt tās paņemtu, un jūs saņemat abus potus.} \\
            }
        }
    }
     \\ \\

    \theutterance \stepcounter{utterance}  
    & & & \multicolumn{2}{p{0.3\linewidth}}{
        \cellcolor[rgb]{0.9,0.9,0.9}{
            \makecell[{{p{\linewidth}}}]{
                \texttt{\tiny{[GM$|$GM]}}
                \texttt{continue} \\
            }
        }
    }
    & & \\ \\

    \theutterance \stepcounter{utterance}  
    & & & \multicolumn{2}{p{0.3\linewidth}}{
        \cellcolor[rgb]{0.9,0.9,0.9}{
            \makecell[{{p{\linewidth}}}]{
                \texttt{\tiny{[GM$|$GM]}}
                \texttt{Sveiki! Lieliski, tas labi saskan. Man savukārt poti (pots) ir pilnīgi bezvērtīgi (0), tāpēc jūs varat paņemt abus potus. Es augstu vērtēju bultskrūvi un knaibles. Soma man ir tikai nedaudz vērtīga. Kā jums ir ar bultskrūvi un somu? Ja jūs tās nevērtējat, es labprāt tās paņemtu, un jūs saņemat abus potus.} \\
            }
        }
    }
    & & \\ \\

    \theutterance \stepcounter{utterance}  
    & & \multicolumn{4}{p{0.6\linewidth}}{
        \cellcolor[rgb]{0.9,0.9,0.9}{
            \makecell[{{p{\linewidth}}}]{
                \texttt{\tiny{[P1$\langle$GM]}}
                \texttt{Sveiki! Lieliski, tas labi saskan. Man savukārt poti (pots) ir pilnīgi bezvērtīgi (0), tāpēc jūs varat paņemt abus potus. Es augstu vērtēju bultskrūvi un knaibles. Soma man ir tikai nedaudz vērtīga. Kā jums ir ar bultskrūvi un somu? Ja jūs tās nevērtējat, es labprāt tās paņemtu, un jūs saņemat abus potus.} \\
            }
        }
    }
    & \\ \\

    \theutterance \stepcounter{utterance}  
    & \multicolumn{4}{p{0.6\linewidth}}{
        \cellcolor[rgb]{0.9,0.9,0.9}{
            \makecell[{{p{\linewidth}}}]{
                \texttt{\tiny{[P1$\rangle$GM]}}
                \texttt{Skan lieliski! Ņemot vērā jūsu vērtības, ierosinu šādu sadalījumu: jūs saņemat bultskrūvi un abas knaibles (ko jūs augstu vērtējat), es saņemu abus potus un somu. Tā kā soma man ir vērtīga (2), bet jums tikai nedaudz, tas maksimizē mūsu kopējo rezultātu. Vai piekrītat?} \\
\\ 
\texttt{{[}Priekšlikums: 2 poti, 1 soma{]}} \\
            }
        }
    }
    & & \\ \\

    \theutterance \stepcounter{utterance}  
    & & & \multicolumn{2}{p{0.3\linewidth}}{
        \cellcolor[rgb]{0.9,0.9,0.9}{
            \makecell[{{p{\linewidth}}}]{
                \texttt{\tiny{[GM$|$GM]}}
                \texttt{proposal} \\
            }
        }
    }
    & & \\ \\

    \theutterance \stepcounter{utterance}  
    & & & \multicolumn{2}{p{0.3\linewidth}}{
        \cellcolor[rgb]{0.9,0.9,0.9}{
            \makecell[{{p{\linewidth}}}]{
                \texttt{\tiny{[GM$|$GM]}}
                [0, 1, 0, 2]
            }
        }
    }
    & & \\ \\

    \theutterance \stepcounter{utterance}  
    & & \multicolumn{4}{p{0.6\linewidth}}{
        \cellcolor[rgb]{0.9,0.9,0.9}{
            \makecell[{{p{\linewidth}}}]{
                \texttt{\tiny{[GM$\rangle$P2]}}
                \texttt{Otrs spēlētājs ir iesniedzis slepenu priekšlikumu. Tagad ir jūsu kārta iesniegt savu priekšlikumu.} \\
            }
        }
    }
    & \\ \\

    \theutterance \stepcounter{utterance}  
    & & & \multicolumn{4}{p{0.6\linewidth}}{
        \cellcolor[rgb]{0.9,0.9,0.9}{
            \makecell[{{p{\linewidth}}}]{
                \texttt{\tiny{[GM$\langle$P2]}}
                \texttt{{[}Priekšlikums: 1 bultskrūve, 1 soma, 2 knaibles{]}} \\
            }
        }
    }
     \\ \\

    \theutterance \stepcounter{utterance}  
    & & & \multicolumn{2}{p{0.3\linewidth}}{
        \cellcolor[rgb]{0.9,0.9,0.9}{
            \makecell[{{p{\linewidth}}}]{
                \texttt{\tiny{[GM$|$GM]}}
                \texttt{proposal} \\
            }
        }
    }
    & & \\ \\

    \theutterance \stepcounter{utterance}  
    & & & \multicolumn{2}{p{0.3\linewidth}}{
        \cellcolor[rgb]{0.9,0.9,0.9}{
            \makecell[{{p{\linewidth}}}]{
                \texttt{\tiny{[GM$|$GM]}}
                [1, 1, 2, 0]
            }
        }
    }
    & & \\ \\

    \theutterance \stepcounter{utterance}  
    & & & \multicolumn{2}{p{0.3\linewidth}}{
        \cellcolor[rgb]{0.9,0.9,0.9}{
            \makecell[{{p{\linewidth}}}]{
                \texttt{\tiny{[GM$|$GM]}}
                [[0, 1, 0, 2], [1, 1, 2, 0]]
            }
        }
    }
    & & \\ \\

\end{supertabular}
}

\end{document}
